\section{Compression and Reference Orbit Reuse Modes}
\label{subsec:ref-orbit-compression-reuse}

Two orthogonal storage decisions are made while iterating:

\begin{enumerate}
  \item \textbf{Orbit storage} for perturbation use (\code{PExtras}):
  \begin{itemize}
    \item \code{Disable}: store every \((\hat{x}_n,\hat{y}_n)\) uncompressed.
    \item \code{SimpleCompression}: store a compressed subset of iterations using an error exponent
          determined by \code{Fractal::CompressionError}.
    \item \code{MaxCompression}: store the Imagina-compatible compressed
          representation used by \code{CompressMax}.
    \item \code{Bad}: store orbit values plus underflow/diagnostic flags.
  \end{itemize}
  \item \textbf{Reuse storage} for intermediate-precision regeneration (\code{ReuseMode}):
  \begin{itemize}
    \item \code{SaveForReuse1/2}: store uncompressed \code{mpf\_t} reuse entries.
    \item \code{SaveForReuse3}: store an intermediate-compressed reuse stream.
    \item \code{SaveForReuse4}: currently unsupported; the implementation
          reaches an assertion instead of producing a reuse stream.
  \end{itemize}
\end{enumerate}

The next two sections describe these mechanisms in more detail.

\subsection{Reuse-based orbit regeneration}
\label{subsec:ref-orbit-reuse}

We consider iteration of the quadratic map
\[
f_c(z) = z^2 + c,\qquad z,c \in \mathbb{C},
\]
and distinguish between an \emph{authoritative} reference orbit computed in very high precision and a hierarchy of
\emph{intermediate precision} reference orbits constructed via perturbation from that authoritative orbit.
The central motivation for this construction is performance at extreme zoom depths: while the authoritative orbit
is expensive to compute, it need only be regenerated infrequently as the view parameter drifts.
Intermediate precision orbits, by contrast, are substantially cheaper to evaluate and can be reused across many
incremental zoom steps before their accuracy envelope is exceeded.

\subsubsection{Authoritative Reference Orbit}

\paragraph{Notation.}
The earlier perturbation sections
(\cref{sec:perturbation-concept,sec:perturb-only})
write the reference orbit and parameter as
\(z_n^\star\) and \(c_\star\).
The following uses the indexed notation
\(z^{(0)}_n\) and \(c_0\) instead, so that the superscript can
distinguish the authoritative level~\(0\) from intermediate
levels~\(1,2,\ldots\) in the reuse hierarchy.
The two conventions describe the same object:
\(z^{(0)}_n \equiv z_n^\star\) and \(c_0 \equiv c_\star\).

Let
\[
c_0 \in \mathbb{C}
\]
denote the authoritative reference parameter.
The corresponding authoritative reference orbit is defined by
\begin{equation}
z^{(0)}_0 = 0, \qquad
z^{(0)}_{n+1} = \bigl(z^{(0)}_n\bigr)^2 + c_0,
\label{eq:authoritative_orbit}
\end{equation}
and is computed at the selected finite precision.  It is authoritative for
subsequent perturbative constructions in the sense that they consume its stored
samples; it is not an exact real orbit.

At very deep zooms, recomputing~\eqref{eq:authoritative_orbit} at every navigation step would dominate runtime,
as the required precision grows with zoom depth.
The reuse-based framework therefore seeks to minimize how often this authoritative orbit must be recalculated,
while still maintaining numerical correctness.

\subsubsection{Perturbation Formulation}

For any nearby parameter
\[
c = c_0 + \Delta c,
\]
the corresponding orbit may be written as
\[
z_n(c) = z^{(0)}_n + \Delta z_n,
\]
where $\Delta z_n$ represents the perturbation relative to the authoritative orbit.
Substituting into the defining recurrence yields the exact perturbation equation
\begin{align}
\Delta z_0 &= 0, \label{eq:perturbation-init} \\
\Delta z_{n+1}
&= 2 z^{(0)}_n \Delta z_n + (\Delta z_n)^2 + \Delta c.
\label{eq:perturbation_exact}
\end{align}
Equation~\eqref{eq:perturbation_exact} is the fundamental relation used both to construct intermediate reference orbits
and to evaluate per-pixel perturbations during rendering.

\subsubsection{Intermediate Precision Reference Orbits}

An intermediate precision reference orbit is defined at a parameter
\[
c_1 = c_0 + \Delta c_1,
\]
where $\Delta c_1$ is chosen such that the perturbation $\Delta z^{(1)}_n$ remains bounded within a fixed target accuracy
(e.g.\ absolute error $\lesssim 10^{-100}$) when represented in a chosen intermediate precision arithmetic.

The intermediate reference orbit is defined by
\[
z^{(1)}_n \equiv z^{(0)}_n + \Delta z^{(1)}_n,
\]
where $\Delta z^{(1)}_n$ is obtained by iterating
\begin{equation}
\Delta z^{(1)}_{n+1}
= 2 z^{(0)}_n \Delta z^{(1)}_n
  + \bigl(\Delta z^{(1)}_n\bigr)^2
  + \Delta c_1
\label{eq:intermediate_perturbation}
\end{equation}
in the intermediate precision format.

The cost of evaluating~\eqref{eq:intermediate_perturbation} is substantially lower than that of recomputing
the authoritative orbit~\eqref{eq:authoritative_orbit}.
As a result, once an intermediate reference orbit has been established, it may be reused across many subsequent zoom
or pan operations, as long as the view parameter remains within its validity radius.
Only when accumulated drift causes the perturbation to exceed the fixed accuracy envelope does a new authoritative
orbit need to be computed.

In the single-thread execution path (beginning at
\code{Add\-Perturbation\-Reference\-Point\-ST}), this construction is performed
sequentially using the cached authoritative
orbit samples $\{z^{(0)}_n\}$ as coefficients.
No approximation beyond finite-precision rounding is introduced; the intermediate orbit is mathematically equivalent
to directly iterating $f_{c_1}$, subject only to the chosen precision bound.

\subsubsection{Per-Pixel Perturbation from an Intermediate Orbit}

For an individual pixel parameter
\[
c_{\text{px}} = c_1 + \delta c,
\]
with $|\delta c| \ll |\Delta c_1|$, the final orbit is expressed as
\[
z_n(c_{\text{px}}) = z^{(1)}_n + \delta z_n,
\]
where $\delta z_n$ satisfies
\begin{equation}
\delta z_{n+1}
= 2 z^{(1)}_n \delta z_n
  + (\delta z_n)^2
  + \delta c.
\label{eq:pixel_perturbation}
\end{equation}
Because $|\delta c|$ is small, $\delta z_n$ remains well within the same fixed precision envelope used for the
intermediate reference orbit.

From a performance perspective, this two-level perturbation hierarchy is
beneficial when interacting with the Mandelbrot. The expensive high-precision
authoritative orbit is amortized over many intermediate orbits, and each
intermediate orbit in turn supports an entire image worth of per-pixel
perturbations. As zoom depth increases, most navigation steps therefore reuse
existing intermediate data, with authoritative recomputation occurring only
sporadically.

The two-level approach nevertheless does have an important downside: instead of
maintaining only a low-precision copy of the reference orbit, the system must also
store the authoritative orbit samples $\{z^{(0)}_n\}$ to support intermediate
reconstruction.  This requirement increases memory usage and data transfer requirements,
potentially impacting performance.

\subsubsection{SaveForReuse Modes}

Both \code{SaveForReuse1} and \code{SaveForReuse2} store the same uncompressed
MPIR reuse-entry format.  They differ in execution strategy.

\paragraph{SaveForReuse1.}
In \code{SaveForReuse1}, the authoritative orbit samples $\{z^{(0)}_n\}$ are stored after initial computation and reused
whenever intermediate reference orbits are constructed.
This approach avoids recomputation of~\eqref{eq:authoritative_orbit} and ensures that all perturbative steps are driven by the same
authoritative data.

\paragraph{SaveForReuse2.}
\code{SaveForReuse2} offloads reuse-entry construction to the reuse worker,
whereas \code{SaveForReuse1} stores entries on the main reference thread.  Both
feed the same later perturbation recurrence.

The two modes have the same intended stored representation and mathematical
recurrence, but the code does not establish a general cross-schedule bitwise
identity guarantee.  \code{SaveForReuse3} instead uses the simple intermediate
MPIR compressor at a fixed extra precision and preserves monotone decompression
semantics.  \code{SaveForReuse4} is broken and must not be presented as an
available maximal intermediate compressor.

\subsection{Reference Compression}
\label{sec:ref-compression-zhuoran}

The idea in this section was developed by Zhuoran
\cite{Zhuoran2023ReferenceCompression} and \FractalShark{} implements it.

\subsubsection{Motivation}
For deep zoom escape-time fractals, pixel evaluation has become extremely fast
(e.g., via perturbation and GPU acceleration), while \emph{reference-orbit}
construction often remains comparatively expensive and can dominate total render
time.  Additionally, when the reference orbit is large, it can consume
significant amounts of memory.  The goal of \emph{reference compression} is to
store this potentially-large reference orbit in a compact form that can be
transmitted or cached, then \emph{reconstructed} efficiently. The implementation
stores a sparse set of \emph{waypoints} and fills gaps by recomputing
intermediate iterations in reduced precision.  Its waypoint thresholds are
implementation criteria, not a formal global error proof.

Two final-orbit formats must be distinguished.  \code{Simple\allowbreak Compression} uses
\code{Runtime\allowbreak Decompressor}: it stores reduced values of type~\code{T}, supports
binary-search anchoring for random access, and maintains two cached states for
sequential access.  \code{Max\allowbreak Compression} uses \code{CompressMax}: it begins in
absolute mode, can switch to perturbation/rebase mode, fully materializes the
orbit during decompression, and then applies backward corrections.  The latter
is therefore not streamable in the same sense as the simple decompressor.

\subsubsection{High-precision reference and low-precision surrogate}
Let the selected reference samples before compression be
\(
z_n^{\mathrm{HP}} \in \mathbb{C}
\)
for \(n=0,\dots,N\).
During compression and decompression, the implementation also maintains a low-precision surrogate
\(
\hat z_n \in \mathbb{C}
\)
computed by iterating the same recurrence in a cheaper numeric type.

Because the dynamics are initially stable to rounding and only later amplify
numerical differences, \(\hat z_n\) typically tracks \(z_n^{\mathrm{HP}}\) for
many iterations before diverging.  This observation enables a streaming scheme:
store occasional reduced anchors (waypoints) and recover omitted states by
recomputing them in low precision between anchors.

\subsubsection{Waypoint selection criterion}
A \emph{waypoint} stores the iteration index and a corrective payload.
In the simple mode, the payload is the stored reduced value of type~\code{T} at
iteration \(n_k\), not an exact high-precision anchor.
Waypoints are chosen by simulating \(\hat z_n\) in low precision while comparing
against \(z_n^{\mathrm{HP}}\).

The implemented simple-compression test is

\begin{equation}
  \left((\Re e_n)^2+(\Im e_n)^2\right)
  \code{CompressionError}
  \;\ge\; (\Re z_n)^2+(\Im z_n)^2,
  \label{eq:relerr-test}
\end{equation}

where $e_n=z_n-\hat z_n$.  Whenever the test succeeds, the compressor emits a
waypoint at \(n\) and resets the replay state to the stored reduced value:

\begin{equation}
  \text{if waypoint at } n:\qquad \hat z_n \leftarrow w_n.
  \label{eq:waypoint-reset}
\end{equation}

Between waypoints, the compressor stores nothing, relying on \(\hat z\)-iteration
to reconstruct the missing values later.

\paragraph{Compressed representation.}
The compressed reference is thus a sparse, ordered list
\(
\mathcal{W} = \{(n_k, w_k)\}_{k=0}^{K-1}
\)
with strictly increasing indices \(0 \le n_0 < n_1 < \cdots < n_{K-1} \le N\),
where the payload \(w_k\) depends on the mode (\cref{sec:refcomp-pert}).

\subsubsection{Decompression by replay}

Decompression replays the same low-precision recurrence and applies waypoints as
hard resets:

\begin{equation}
  \hat z_{n+1} = \hat z_n^2 + c \quad\text{(low precision)},\qquad
  \text{and if } n = n_k:\ \hat z_n \leftarrow w_k.
  \label{eq:decompress-replay}
\end{equation}

If the decompressor uses the same reduced arithmetic, recurrence, and waypoint
indices, its replay matches the compressor's replay state.  The placement test
controls when resets occur but does not, by itself, prove a global error bound.

\paragraph{On-the-fly use during rendering.}
Because decompression is a forward scan with occasional resets, it can be
performed \emph{streaming} (no random access required), and can be interleaved
with perturbation-based pixel evaluation by iterating \(\hat z_n\) alongside the
main render loop.

\subsubsection{Perturbation-assisted compression}

\label{sec:refcomp-pert}
This subsection describes \code{MaxCompression}, not
\code{SimpleCompression}.  \code{CompressMax} uses the Chebyshev norm (the
maximum absolute component), begins in absolute mode, and switches only after
the orbit norm falls below $1/16$.  In perturbation mode it tracks a reference
index~$j$ and emits a waypoint when $j$ reaches the previous waypoint or when
the configured error test requires one.

Waypointing only the absolute states \(z_n\) is useful but may not achieve high
compression at extreme depths, where the orbit's sensitivity forces frequent
anchors.  Compression improves substantially by switching to a perturbation form
once the orbit becomes suitable.

\paragraph{Self-referenced perturbation.}
Choose a \emph{base} (reference) iteration \(r\) and represent subsequent states as

\begin{equation}
  z_n = z_r + \Delta z_n, \qquad n \ge r,
  \label{eq:self-ref-perturbation}
\end{equation}

where \(z_r\) is treated as a reference and \(\Delta z_n\) is (ideally) small.
In standard perturbation rendering, \(\Delta z\) is evolved using derivatives
around a known reference orbit.  For \emph{reference compression}, the only
available reference is the reference itself, so the scheme starts in the
absolute waypoint mode (previous sections) and transitions to perturbation only
when \(\lVert \Delta z \rVert\) is sufficiently small that perturbation yields
meaningful extra precision in the reduced type.

\paragraph{Switch criterion.}
When the implemented Chebyshev norm satisfies
$\max(|\Re z_n|,|\Im z_n|)<1/16$, the compressor may begin emitting waypoints
that store \(\Delta z_n\) rather than \(z_n\).  After the switch, a waypoint
payload becomes

\begin{equation}
  w_k = \Delta z_{n_k} \quad\text{(instead of } z_{n_k}\text{)}.
  \label{eq:pert-waypoint-payload}
\end{equation}

\paragraph{Rebasing synchronization.}
Perturbation implementations commonly \emph{rebase} to prevent \(\Delta z\) from
growing too large: at some iteration \(b\), the reference is updated so that
\(\Delta z_b \leftarrow 0\) and the base index becomes \(r \leftarrow b\).
During decompression, if rebases occur at different iterations than during
compression, the decompressor will interpret stored \(\Delta z\) against the
wrong base \(z_r\), producing catastrophic corruption.

A robust compressed format therefore must encode enough information to keep
rebases in lockstep.  One approach is to include a single \emph{rebase flag} per
waypoint indicating whether a rebase occurs at that iteration; if a rebase
occurs at an iteration where no waypoint would otherwise be required, the
compressor must still record the iteration (e.g., by forcing a waypoint or
storing an index) so the decompressor can rebase at the same \(b\).

\subsubsection{Error propagation and correction}

\code{CompressMax} decompression first materializes the full reconstructed orbit.
The implementation then calls \code{CorrectOrbit} over waypoint-delimited
segments.  That routine walks backward while accumulating the complex derivative
of the iterated map, forms the complex difference between the stored waypoint
and the reconstructed endpoint, divides that difference by the accumulated
derivative, and applies the resulting correction to the earlier sample.  This
is the implemented backward correction; the code is not a literal recurrence
$e_n=e_{n+1}/(2z_n)$ at every sample.  Because correction requires later
waypoints and materialized earlier samples, maximal decompression cannot provide
the simple format's one-pass random/sequential runtime interface.

\subsubsection{Algorithm sketch}

The following high-level procedure summarizes the method.

\begin{enumerate}
  \item \textbf{Compression pass:}
    \begin{enumerate}
      \item Initialize low-precision \(\hat z_0 \leftarrow 0\); iterate forward.
      \item At each \(n\), evaluate relative error \eqref{eq:relerr-test}.
      \item If error too large, emit waypoint \((n, z_n^{\mathrm{HP}})\) (or \((n,\Delta z_n)\) after perturbation switch),
            set \(\hat z_n \leftarrow z_n^{\mathrm{HP}}\) (or reset perturbation state), and continue.
      \item If in perturbation mode, record/synchronize any rebase events (bit flag or forced waypoint).
      \item Insert additional waypoints before an uncorrected state would be
            reused as a perturbation base.
    \end{enumerate}

  \item \textbf{Decompression pass:}
    \begin{enumerate}
      \item Replay the same low-precision recurrence, applying waypoints as resets.
      \item In perturbation mode, rebase exactly when signaled by the compressed stream.
      \item After materializing the orbit, call \code{CorrectOrbit} to apply
            derivative-based backward corrections over waypoint-delimited segments.
    \end{enumerate}
\end{enumerate}

\subsubsection{Practical tradeoffs}

Reference compression exposes a clear accuracy--size--time trade:
\begin{itemize}
  \item Smaller tolerance \(\varepsilon\) \(\Rightarrow\) more waypoints \(\Rightarrow\) larger compressed size but tighter fidelity.
  \item Larger gaps reduce storage but increase decompression work and may require more frequent correction points,
        especially in perturbation mode.
  \item Perturbation-mode waypoints (\(\Delta z\)) can greatly improve compression when \(\Delta z\) remains small for long
        stretches, but require careful rebase synchronization and correction to avoid catastrophic desynchronization.
\end{itemize}

Overall, the simple format provides compact replay with random and sequential
access, while the maximal format trades full materialization and backward
correction for a different compression ratio.  They should not be treated as
one interchangeable streaming representation.

\subsubsection{Supporting both random and sequential access patterns}
\label{sec:refcomp-runtime-access}

Reference-orbit consumers exhibit two different access patterns.
On the one hand, perturbation-based evaluation and on-the-fly decompression
naturally traverse the orbit sequentially, requesting
\(z_0, z_1, z_2, \dots\) in order.  On the other hand, higher-level algorithms
such as linear (or affine) approximation, reuse heuristics, and error analysis
require \emph{direct access} to arbitrary iteration indices in order to evaluate
local behavior around a chosen anchor.  A decompression strategy optimized only
for sequential replay performs poorly for such random probes, while a strategy
optimized only for random access would impose unnecessary overhead on the
dominant sequential case.

To accommodate both use cases efficiently, the runtime decompressor implements a
hybrid access scheme: binary-search-based anchoring for random access, combined
with cached linear scans for sequential access.

\paragraph{Random access via waypoint search.}
The compressed reference orbit is stored as a strictly increasing sequence of
waypoints
\(
\{(i_k, z_{i_k})\}_{k=0}^{K-1}
\),
where \(i_k\) denotes the uncompressed iteration index and \(z_{i_k}\) the saved
state.  Given a request for an arbitrary iteration \(n\), the decompressor locates the nearest
preceding waypoint
\begin{equation}
  k^\star = \max\{k \mid i_k \le n\},
  \label{eq:waypoint-search}
\end{equation}
using a binary search over the waypoint indices in
\(\mathcal{O}(\log K)\) time.  Reconstruction then proceeds by replaying the
low-precision recurrence forward from that anchor for
\(\Delta = n - i_{k^\star}\) steps:
\begin{equation}
  z_{t+1} \leftarrow z_t^2 + c,
  \qquad t = i_{k^\star}, \dots, n-1,
  \label{eq:replay-recurrence}
\end{equation}
with the same reduction/normalization operations used during compression.  This
provides efficient, bounded-cost random access suitable for linear approximation
and analysis routines that must probe isolated iteration indices.

\paragraph{Sequential access via cached linear scans.}
In contrast, perturbation pipelines typically consume the reference orbit in
monotone order.  Repeating a binary search for every \(n+1\) would be wasteful,
so the decompressor maintains a small cache of recently reconstructed states,
including both their uncompressed and compressed indices.  When a request
targets the immediate successor of a cached iteration, the decompressor advances
the state in-place:
\begin{itemize}
  \item If the next uncompressed index lies strictly between two waypoints, a
        single recurrence step is performed.
  \item If the next index coincides with a waypoint, the cached state is replaced
        by the stored waypoint value and the compressed index is advanced.
\end{itemize}
A dual-entry cache tolerates minor access jitter by retaining both the most
recent and the previously-recent state; exact cache hits are returned
immediately, and near-sequential requests avoid both binary search and replay
from older anchors.

\paragraph{Combined complexity and practical behavior.}
This hybrid design yields the desirable properties of both approaches:
\begin{align}
  \text{sequential access} &: \quad \mathcal{O}(1)\ \text{amortized per iteration}, \label{eq:sequential-complexity} \\
  \text{random access} &: \quad \mathcal{O}(\log K + \Delta), \label{eq:random-access-complexity}
\end{align}
where \(\Delta\) is the distance from the nearest preceding waypoint.  In typical
compressed references, \(\Delta\) is small relative to the full orbit length,
while sequential traversal dominates overall access volume.  The result is a
single runtime decompressor that efficiently supports both perturbation-driven
sequential replay and analysis-driven random probing, without duplicating
reference representations or specializing the data structure to a single access
pattern.

\paragraph{Example usage.}
See \cref{fig:view27-1-quadrillion} for an example of a deep-zoom render
that requires reference compression to complete.

\begin{figure}[!htbp]
  \centering
  \includegraphics[width=0.9\linewidth]{view27-1-quadrillion.png}
  \caption[View \#27 at 1 quadrillion iterations]{View \#27, a period ~28-billion point rendered at 1 quadrillion
  iterations per pixel, requires reference compression to render on most
  hardware because of its memory requirements.  This render takes ~6h with an
  RTX 4090 and AMD 5950X and 128GB RAM with the multithreaded reference orbit
  backend and is effectively the hardest point \FractalShark{} has likely ever
  rendered.}
  \label{fig:view27-1-quadrillion}
\end{figure}
